% ============================================================================
\clearpage
\section*{v4.9.0-r29 Source-Only: Where the Tachyon Reservoir Could Come From}
\addcontentsline{toc}{section}{v4.9.0-r29 Source-Only: Where the Tachyon Reservoir Could Come From}
\label{sec:v490-r29-reservoir-source}
% ============================================================================

\paragraph{Purpose.}
r28 found that the field stabilises between the pure state and the photon only in exchange with a tachyon reservoir that keeps
being replenished.  On 2026-10-22 the author asked where that reservoir could come from.  This layer tests six candidates that the
document itself offers.  Candidates and criteria were fixed in advance in \texttt{llm/k66\_tachyon\_reservoir\_quelle.py}.  No new
prediction is made, and nothing is compared with data.

\noindent\textbf{Source-only status.}
This v4.9.0-r29 layer is source-only.  It does not build a PDF, does not assign a new DOI,
and does not authorize a public release.  The v4.9.0 Version DOI
\href{https://doi.org/10.5281/zenodo.23534954}{10.5281/zenodo.23534954} and the Concept DOI
\href{https://doi.org/10.5281/zenodo.19210728}{10.5281/zenodo.19210728} are unchanged by it.

\subsection*{v4.9.0-r29 criteria and candidates}
\label{subsec:v490-r29-criteria}

A source must (i) supply fresh tachyons again and again, (ii) supply them in the tachyon reference state, polarised on the
partner-port axis, \(b_T=(-1,0)\), (iii) exchange with the system coherently, and (iv) imply no forbidden statement.  The candidates
are: (K1) pairs at a horizon (Theorems~\ref{thm:hawking-triolic}, \ref{thm:core}); (K2) virtual pairs from the vacuum
(Proposition~\ref{prop:shared-vacuum-planck}); (K3) the cascade flow of Theorem~\ref{thm:dark-energy}, \(P(S^8\to S^4)=\bsym^2\);
(K4) the other modes of the field; (K5) a thermal bath through the Wick rotation (Theorem~\ref{thm:wick}); (K6) the tachyon hit
chain of r06.

\subsection*{v4.9.0-r29 results}
\label{subsec:v490-r29-results}

\paragraph{Theoremlet v4.9.0-r29.1.}
\begin{enumerate}[leftmargin=2em]
\item[(a)] \emph{The vacuum supplies tachyons in the reference state.}  By r24 the vacuum \(\zminus=0\) of the main body is the reference
frame R, the pure state on the compensator axis.  Its image under the tachyon map \(z\mapsto iz\) (r28) is exactly \(b_T=(-1,0)\).  A
matter--tachyon pair created from the vacuum therefore hands over a tachyon in its reference state.  K2 meets (ii).
\item[(b)] \emph{Horizons do not.}  The horizon \(\beta\to1\) is the photon fixed point of Convention F (r23).  There \(b_T=0\), and the
attractor of r28.2(c) is \(|b|=0\): no stabilisation.  In Theorem~\ref{thm:hawking-triolic} the tachyon also falls inwards, away
from an outside system.  K1 fails (ii).
\item[(c)] \emph{The other modes and a warm bath only partly.}  Other modes dephase (r28.2(b)) and carry no polarisation, so the attractor
is \(0\) (K4 fails (ii)).  A thermal bath has \(|b_T|<1\), and the attractor falls with it: \(|b|=0.500\), \(0.193\), \(0\) for
\(|b_T|=1\), \(0.5\), \(0\) at \(\sin^2\theta=1/3\).  The zero-temperature limit is the vacuum of (a) (K5 reduces to K2).
\item[(d)] \emph{The hit chain distributes, it does not replenish.}  One tachyon meets many matter systems at once (r06).  After the
first meeting it is no longer in its reference state.  The pair event at the point of infinite speed (r06) is a pair creation,
that is K2.
\item[(e)] \emph{The cascade fixes the strength, if the exchange is coherent.}  With the exchange strength \(\sin^2\theta=\bsym^2\), the
cascade probability, the coherent exchange of r28 has the self-consistent attractor \(|b|=1/2\) exactly at the Hopf point.  Read
instead as a Markov replacement, \(\rho\mapsto(1-p)\rho+p\rho_T\) with \(p=\bsym^2\), the attractor is \(0.378\) and lies off the axis
(\(41^\circ\)).  It is self-consistent only at \(p=0.3865\), which corresponds to no point named in the document.
\end{enumerate}

\paragraph{Proof sketch.}
(a) \(\sigma_y\rho_R\sigma_y\) with \(\rho_R=\tfrac12(1+\sigma_x)\).  (b), (c) iteration of rotation \(\pi/6\) plus exchange to the fixed point,
for the stated \(b_T\).  (e) the same for the coherent partial swap and for the Markov replacement; root of the self-consistency
condition by bisection.  \(\square\)

\paragraph{Bookkeeping.}
At the attractor each kick leaves the incoming tachyon changed, from \(b_T=(-1,0)\) to \((-0.750,0.144)\).  The reservoir takes up
polarisation and must be replaced.  The exported quarter of r11 goes there.  A steady state therefore needs one fresh vacuum pair
per kick.  The rate of pair creation is not fixed by anything in the present layers.

\paragraph{Reading.}
Of the six candidates, only the vacuum supplies tachyons in the right state.  The proposed source is the pair creation from the
vacuum: the matter partner of each pair continues, and the tachyon partner serves once as the measuring device and the reservoir,
coupling through time reversal (r28.1).  The cascade fits the exchange strength \(\bsym^2\), provided the exchange is coherent and
not a classical replacement.  In this picture the matter system holds its state because the vacuum keeps supplying fresh
partners.  Three points stay open: the rate of pair creation, why the exchange should be coherent, and where the absorbed
polarisation of the used tachyons goes.

\subsection*{v4.9.0-r29 status in the three frames}
\label{subsec:v490-r29-frames}

\begin{center}
\small
\begin{tabular}{@{}>{\raggedright\arraybackslash}p{0.30\textwidth}>{\raggedright\arraybackslash}p{0.62\textwidth}@{}}
\toprule
result & M / T / P\\
\midrule
r29.1(a) vacuum pairs & the vacuum is R (pure); the pair has its matter member in M and its tachyon member in T, in the reference state\\
r29.1(b) horizon & lies at P, where the reservoir is unpolarised\\
r29.1(c), (d) & frame-free statements about polarisation; P gives \(0\)\\
r29.1(e) cascade strength & \(\bsym^2\) is a matter-frame quantity (r22); the self-consistency holds at the Hopf point only\\
\bottomrule
\end{tabular}
\end{center}

\subsection*{v4.9.0-r29 register entries}
\label{subsec:v490-r29-register}

\begin{center}
\small
\begin{tabular}{@{}>{\raggedright\arraybackslash}p{0.34\textwidth}>{\raggedright\arraybackslash}p{0.58\textwidth}@{}}
\toprule
statement & status after v4.9.0-r29\\
\midrule
source of the tachyon reservoir & proposed: pair creation from the vacuum R (r29.1(a)); horizon, other modes, hit chain excluded; warm bath only partial\\
exchange strength & \(\bsym^2\) from the cascade, if coherent; Markov reading not self-consistent at a named point (r29.1(e))\\
open & rate of pair creation, coherence of the exchange, fate of the used tachyons\\
\bottomrule
\end{tabular}
\end{center}

\subsection*{v4.9.0-r29 ledger and guard update}
\label{subsec:v490-r29-ledger-guard}

\[
R_{\rm res}^{v490\text{-}r29}=(0,0,1)
\]
for (candidates tested, vacuum pairs proposed, rate and coherence open).

The layer spends no guard:
\[
R_{\rm glob}^{(v490\text{-}r29,+)}=(0,1,1,1,1),
\qquad
\Gamma_{\rm guard}^{v490\text{-}r29}=(1,1,1,1,1),
\qquad
\beta_4=1/256\ \text{unspent}.
\]
It does not imply exact finite \(C5\) zero closure, exact all-readout terminality,
\(C6/T2\), a global photon-clock/RH theorem, or global all-readout closure.

